\documentclass[../../main.tex]{subfiles}
\graphicspath{{\subfix{./image/}}{\subfix{../../image/}}}
\begin{document}
% SOURCE p=196; printed=181; object=ch06-basic-title; region=middle
\section{$L^p$ 空间的基本理论}\label{section:ch06-lp-basic}
% SOURCE p=196; printed=181; object=ch06-fixed-measure-space; region=middle
本章始终在一个固定的测度空间 $(X,\mathscr M,\mu)$ 上讨论。
% SOURCE p=196; printed=181; object=ch06-lp-definition; region=lower
\begin{definition}[$L^p$ 空间]\label{definition:ch06-lp}
若 $f$ 是 $X$ 上的可测函数，且 $0<p<\infty$，定义
\begin{align*}
\|f\|_p=\left[\int|f|^p\,\mathrm d\mu\right]^{1/p},
\end{align*}
允许 $\|f\|_p=\infty$。定义
\begin{align*}
L^p(X,\mathscr M,\mu)=\{f:X\to\mathbb C:f\text{ 可测且 }\|f\|_p<\infty\}.
\end{align*}
不致混淆时，简记为 $L^p(\mu)$、$L^p(X)$ 或 $L^p$。与 $L^1$ 的情形一样，两个几乎处处相等的函数视为 $L^p$ 中同一个元素。
\end{definition}
% SOURCE p=196; printed=181; object=ch06-little-lp-definition; region=lower
\begin{definition}[$\ell^p$ 空间]\label{definition:ch06-little-lp}
若 $A$ 为非空集合，令 $\mu$ 为 $(A,\mathscr P(A))$ 上的计数测度，定义 $\ell^p(A)=L^p(\mu)$；将 $\ell^p(\mathbb N)$ 简记为 $\ell^p$。
\end{definition}
% SUPPLEMENT: id=v14-ch06-lp-definition-details; parent=ch06-lp-definition; kind=proof+motivation+hypotheses; origin=independent-derivation
\begin{proof}
\textbf{等价类的良定义性（重构新增）。}可测函数之间的几乎处处相等是等价关系：自反与对称直接成立，传递性只需把两个可测例外零集取并。若 $f=f'$、$g=g'$ 几乎处处，则在这两个零集之并以外，$af+bg=af'+bg'$，而 $|f|^p=|f'|^p$。积分忽略可测零集，故 $\|f\|_p=\|f'\|_p$，线性运算也不依赖代表。计数测度唯一的零集是空集，因此在 $\ell^p(A)$ 中，几乎处处相等就是处处相等。
\end{proof}
\begin{remark}
\textbf{动机与条件说明（重构新增）。}这两个定义把函数的大小汇总为一个数，计数测度时汇总即为 $\sum_{a\in A}|f(a)|^p$；任意指标集上的非负和取有限子和的上确界。$p>0$ 保证实数幂在非负数上有意义并保持大小关系，有限 $p$ 暂不包含稍后的本质上确界端点。可测性使被积函数可测，有限积分选出具有有限大小的元素；把几乎处处相等者识别，才能使“大小为零”意味着零元素。值域为 $\mathbb C$ 给出复向量空间，改成 $\mathbb R$ 可得到相同论证的实版本。定义不要求测度有限、$\sigma$ 有限或完备；所用例外集均为可测零集。非空 $A$ 沿原书约定，若允许 $A=\varnothing$，所得空间只是零空间，后面的包含命题也涵盖它。$0<p<1$ 时此量通常不是范数，定义并未预先声称三角不等式。
\end{remark}
% END SUPPLEMENT: id=v14-ch06-lp-definition-details
% SOURCE p=196,197; printed=181,182; object=ch06-lp-linear-space; region=bottom-top
\begin{proposition}\label{proposition:ch06-lp-linear-space}
$L^p$ 是向量空间。
\end{proposition}
\begin{proof}
若 $f,g\in L^p$，则
\begin{align*}
|f+g|^p\leqslant[2\max(|f|,|g|)]^p\leqslant2^p(|f|^p+|g|^p),
\end{align*}
所以 $f+g\in L^p$。
\end{proof}
% SOURCE p=197; printed=182; object=ch06-lp-norm-basic-properties; region=top
\begin{proposition}\label{proposition:ch06-lp-basic-properties}
$\|f\|_p=0$ 当且仅当 $f=0$ 几乎处处，且 $\|cf\|_p=|c|\|f\|_p$。
\end{proposition}
% SOURCE p=197; printed=182; object=ch06-lp-norm-discussion; region=top
我们的记号暗示 $\|\cdot\|_p$ 是 $L^p$ 上的范数。事实上，上面的两条性质是显然的，因此唯一的问题是三角不等式。它恰在 $p\geqslant1$ 时成立，所以我们几乎只关注这一情形。
% SUPPLEMENT: id=v14-ch06-lp-linear-details; parent=ch06-lp-linear-space; kind=proof+motivation+hypotheses; origin=independent-derivation
\begin{proof}
\textbf{线性结构及零值性质补写（重构新增）。}零函数可测且积分为零；标量乘法保持可测性，并给 $\int|cf|^p=|c|^p\int|f|^p$，对 $L^p$ 元素各量有限，因此包括 $c=0$ 的情形也合法。结合原证明的加法封闭性和代表无关性，得到向量空间。若 $\int|f|^p=0$，对每个 $n\geqslant1$，集合 $E_n=\{|f|^p>1/n\}$ 满足 $\mu(E_n)/n\leqslant\int|f|^p=0$；因此 $\{f\ne0\}=\bigcup_nE_n$ 为零集。反向由零集不影响积分得到。标量齐次性由上式取 $p$ 次方根得到。
\end{proof}
\begin{remark}
\textbf{条件与范围（重构新增）。}\cref{proposition:ch06-lp-linear-space,proposition:ch06-lp-basic-properties}只用 $p>0$，不需要三角不等式；粗略的常数 $2^p$ 已足以保证加法封闭。有限积分排除了 $0\cdot\infty$ 等不必要歧义。原文“恰在”应理解为对一般测度空间成立的统一理论；某些退化空间在 $p<1$ 时仍可能得到范数，例如单点计数测度下 $\|f\|_p=|f|$。下面原例明确给出三角不等式失败所需的两个不交正有限测度集合。
\end{remark}
% END SUPPLEMENT: id=v14-ch06-lp-linear-details
% SOURCE p=197; printed=182; object=ch06-small-p-counterexample; region=upper
先说明 $p<1$ 时三角不等式为何失败。
\begin{example}\label{example:ch06-small-p-triangle}
若 $a>0$、$b>0$ 且 $0<p<1$，则 $a^p+b^p>(a+b)^p$。因此，只要测度空间中存在两个不交的正有限测度集合，$\|\cdot\|_p$ 就不满足三角不等式。
\end{example}
\begin{proof}
对 $t>0$，有 $t^{p-1}>(a+t)^{p-1}$；从 $0$ 到 $b$ 积分得到 $a^p+b^p>(a+b)^p$。因此，若 $E,F$ 是 $X$ 中不交的正有限测度集合，令 $a=\mu(E)^{1/p}$、$b=\mu(F)^{1/p}$，则
\begin{align*}
\|\chi_E+\chi_F\|_p=(a^p+b^p)^{1/p}>a+b=\|\chi_E\|_p+\|\chi_F\|_p.
\end{align*}
\end{proof}
% SUPPLEMENT: id=v14-ch06-small-p-details; parent=ch06-small-p-counterexample; kind=proof-detail+motivation+hypotheses; origin=independent-derivation
\begin{proof}
\textbf{积分与反例核验（重构新增）。}$p-1<0$ 使幂函数严格递减。两侧在 $(0,b)$ 都可积，第一侧的瑕积分为 $b^p/p$，第二侧为 $((a+b)^p-a^p)/p$；严格正的连续差在任一内部闭子区间上积分为正，故积分不等式严格。乘 $p>0$ 即得原式。$E,F$ 不交使 $|\chi_E+\chi_F|^p=\chi_E+\chi_F$；正测度给 $a,b>0$，有限测度使两示性函数属于 $L^p$。因此确为同一空间中两个元素违反三角不等式。
\end{proof}
\begin{remark}
\textbf{动机与条件说明（重构新增）。}小于一的幂先压缩再取根，会放大分散在不交区域上的合成大小。两个正有限测度块排除了上一注记所说的退化情形；$p=1$ 时严格不等式变成等式。原书这一例不是说每一种测度空间在 $p<1$ 都必定失败。
\end{remark}
% END SUPPLEMENT: id=v14-ch06-small-p-details
% SOURCE p=197; printed=182; object=ch06-holder-introduction; region=middle
$L^p$ 空间理论的基石是 Hölder 不等式，下面推导它。
% SOURCE p=197; printed=182; object=ch06-weighted-arithmetic-geometric; number=6.1; region=middle
\begin{lemma}\label{lemma:ch06-weighted-arithmetic-geometric}
若 $a\geqslant0$、$b\geqslant0$ 且 $0<\lambda<1$，则
\begin{align*}
a^\lambda b^{1-\lambda}\leqslant\lambda a+(1-\lambda)b,
\end{align*}
等号成立当且仅当 $a=b$。
\end{lemma}
\begin{proof}
若 $b=0$，结论显然。否则，两边除以 $b$ 并令 $t=a/b$，归结为证明 $t^\lambda\leqslant\lambda t+(1-\lambda)$，且等号当且仅当 $t=1$。由初等微积分，$t^\lambda-\lambda t$ 在 $t<1$ 时严格递增，在 $t>1$ 时严格递减，所以其最大值 $1-\lambda$ 在 $t=1$ 处取得。
\end{proof}
% SUPPLEMENT: id=v14-ch06-weighted-amgm-details; parent=ch06-weighted-arithmetic-geometric; kind=proof-detail+motivation+hypotheses; origin=independent-derivation
\begin{proof}
\textbf{端点与导数细节（重构新增）。}$b=0$ 时左边为零，右边为 $\lambda a$，故等号恰当 $a=0$。$b>0$ 时 $t\geqslant0$；对 $t>0$，导数为 $\lambda(t^{\lambda-1}-1)$，在 $(0,1)$ 正、在 $(1,\infty)$ 负。函数在 $t=0$ 连续且值为零，严格小于 $1-\lambda$，所以零端点不另产生等号。
\end{proof}
\begin{remark}
\textbf{动机与条件说明（重构新增）。}此不等式把乘积改为可分别积分的加权和。非负性使实幂有定义，两个权重严格介于零和一之间使两项都参与并给严格等号判据。允许 $\lambda=0$ 或 $1$ 时可另约定零次幂并得到平凡边界式，但“等号仅当 $a=b$”便不能保持。取 $\lambda=1/p$ 将使乘积正好变为 $|fg|$。
\end{remark}
% END SUPPLEMENT: id=v14-ch06-weighted-amgm-details
% SOURCE p=197; printed=182; object=ch06-holder; number=6.2; region=middle-lower
\begin{theorem}[Hölder 不等式]\label{theorem:ch06-holder}
设 $1<p<\infty$ 且 $p^{-1}+q^{-1}=1$，即 $q=p/(p-1)$。若 $f,g$ 是 $X$ 上的可测函数，则
\begin{align}
\|fg\|_1\leqslant\|f\|_p\|g\|_q.\label{align:ch06-holder}
\end{align}
特别地，若 $f\in L^p$ 且 $g\in L^q$，则 $fg\in L^1$；在此情形，\eqref{align:ch06-holder}中等号成立当且仅当存在常数 $\alpha,\beta$，其中 $(\alpha,\beta)\ne(0,0)$，使 $\alpha|f|^p=\beta|g|^q$ 几乎处处成立。
\end{theorem}
% SOURCE p=197; printed=182; object=ch06-holder-proof; region=lower
\begin{proof}
若 $\|f\|_p=0$ 或 $\|g\|_q=0$，则 $f=0$ 或 $g=0$ 几乎处处，结论平凡；若 $\|f\|_p=\infty$ 或 $\|g\|_q=\infty$，结论也平凡。此外，若\eqref{align:ch06-holder}对某一对 $f,g$ 成立，则对它们的任意标量倍也成立，因为将 $f,g$ 换成 $af,bg$ 时两边都乘以因子 $|ab|$。所以只需证明 $\|f\|_p=\|g\|_q=1$ 时不等式成立，且等号当且仅当 $|f|^p=|g|^q$ 几乎处处成立。为此，对\cref{lemma:ch06-weighted-arithmetic-geometric}取 $a=|f(x)|^p$、$b=|g(x)|^q$ 及 $\lambda=p^{-1}$，得到
\begin{align}
|f(x)g(x)|\leqslant p^{-1}|f(x)|^p+q^{-1}|g(x)|^q.\label{align:ch06-holder-pointwise}
\end{align}
两边积分给出
\begin{align*}
\|fg\|_1\leqslant p^{-1}\int|f|^p+q^{-1}\int|g|^q
=p^{-1}+q^{-1}=1=\|f\|_p\|g\|_q.
\end{align*}
此处等号成立，当且仅当\eqref{align:ch06-holder-pointwise}几乎处处取等号；由该引理，这又恰好等价于 $|f|^p=|g|^q$ 几乎处处成立。
\end{proof}
% SUPPLEMENT: id=v14-ch06-holder-details; parent=ch06-holder; kind=proof-detail+motivation+hypotheses; origin=independent-derivation
\begin{proof}
\textbf{归一化与全部等号情形（重构新增）。}零范数情形先处理，因此在剩余情形中两范数均正；其中任一个为无穷时右端为无穷，不需要作除法。两者正且有限时，应用原证明于 $F=f/\|f\|_p$、$G=g/\|g\|_q$，再乘两范数即可。原归一化不等式右边与左边之差为非负可积函数，它的积分为零当且仅当它几乎处处为零，这说明积分取等号确能推出逐点取等号。

现仅讨论 $f\in L^p,g\in L^q$。两范数正时，等号等价于 $|f|^p/\|f\|_p^p=|g|^q/\|g\|_q^q$ 几乎处处；这给出所需非零常数对。反之，若 $\alpha|f|^p=\beta|g|^q$，两范数又均正，积分得 $\alpha\|f\|_p^p=\beta\|g\|_q^q$，从而 $\alpha,\beta$ 均非零并恢复上述归一化等式。若 $f=0$ 几乎处处，取 $(\alpha,\beta)=(1,0)$；若 $g=0$ 几乎处处，取 $(0,1)$。反方向中若某个常数为零，另一个非零即强制对应函数为零。因此退化情形也恰好包含在原判据内。
\end{proof}
\begin{remark}
\textbf{动机与逐条件说明（重构新增）。}共轭关系使 $p\lambda=q(1-\lambda)=1$，从而加权几何平均成为 $|fg|$；归一化消去两个独立尺度。$1<p<\infty$ 使权重严格在 $(0,1)$，并使 $q$ 同样有限大于一；端点另在下面讨论。可测性保证积分合法，等号判据的有限积分假设使归一化和非负差积分可用，不能把双方皆无穷的“等号”也按此分类。零范数先处理；一般扩展积分不等式中的零乘无穷依测度论约定取零，且此时乘积确实几乎处处为零。复值函数的相位不影响本式，因为比较的是 $\int|fg|$。整个论证无须 $\mu$ 有限或 $\sigma$ 有限。
\end{remark}
% END SUPPLEMENT: id=v14-ch06-holder-details
% SOURCE p=198; printed=183; object=ch06-conjugate-exponent; region=top
\begin{definition}[共轭指数]\label{definition:ch06-conjugate-exponent}
Hölder 不等式中的关系 $p^{-1}+q^{-1}=1$ 在 $L^p$ 理论中经常出现。若 $1<p<\infty$，则满足这一关系的数 $q=p/(p-1)$ 称为 $p$ 的\textbf{共轭指数}。
\end{definition}
% SUPPLEMENT: id=v14-ch06-conjugate-exponent-remark; parent=ch06-conjugate-exponent; kind=motivation+hypotheses; origin=independent-derivation
\begin{remark}
\textbf{动机与条件说明（重构新增）。}“共轭”在这里记录可相乘后积分的互补幂次，而非复共轭。$p>1$ 保证分母正，$p<\infty$ 给 $q>1$；由 $1/q=1-1/p$ 也可恢复 $p$，故关系对称。定义选取倒数和为一，是为了正好适配 Hölder 估计及其后的对偶配对。
\end{remark}
% END SUPPLEMENT: id=v14-ch06-conjugate-exponent-remark
% SOURCE p=198; printed=183; object=ch06-minkowski; number=6.5; region=upper
\begin{theorem}[Minkowski 不等式]\label{theorem:ch06-minkowski}
若 $1\leqslant p<\infty$ 且 $f,g\in L^p$，则
\begin{align*}
\|f+g\|_p\leqslant\|f\|_p+\|g\|_p.
\end{align*}
\end{theorem}
\begin{proof}
若 $p=1$ 或 $f+g=0$ 几乎处处，结论显然。否则，观察到
\begin{align*}
|f+g|^p\leqslant(|f|+|g|)|f+g|^{p-1},
\end{align*}
然后应用 Hölder 不等式。若 $q$ 是 $p$ 的共轭指数，则 $(p-1)q=p$，因此
\begin{align*}
\int|f+g|^p
&\leqslant\|f\|_p\bigl\||f+g|^{p-1}\bigr\|_q
+\|g\|_p\bigl\||f+g|^{p-1}\bigr\|_q\\
&=(\|f\|_p+\|g\|_p)\left(\int|f+g|^p\right)^{1/q}.
\end{align*}
所以
\begin{align*}
\|f+g\|_p=\left[\int|f+g|^p\right]^{1-1/q}\leqslant\|f\|_p+\|g\|_p.
\end{align*}
\end{proof}
% SUPPLEMENT: id=v14-ch06-minkowski-details; parent=ch06-minkowski; kind=proof-detail+motivation+hypotheses; origin=independent-derivation
\begin{proof}
\textbf{端点与除法合法性（重构新增）。}$p=1$ 时，将复数三角不等式 $|f+g|\leqslant|f|+|g|$ 积分即得结论。$f+g=0$ 几乎处处时左边为零。其余情形中，前面的向量空间命题已在不依赖 Minkowski 的情况下证明 $f+g\in L^p$，所以 $A=\int|f+g|^p$ 正且有限；又 $\int(|f+g|^{p-1})^q=A$，故 Hölder 的两次应用合法。由 $A\leqslant(\|f\|_p+\|g\|_p)A^{1/q}$ 除以正有限的 $A^{1/q}$，再用 $1-1/q=1/p$ 即得原式。
\end{proof}
\begin{remark}
\textbf{动机与条件说明（重构新增）。}三角不等式把积分定义的大小转成可做线性分析的范数。证明把一个 $p$ 次幂拆为一次的 $f$ 或 $g$ 与 $p-1$ 次幂的公共因子，Hölder 的共轭指数恰好重新产生原积分。$f,g\in L^p$ 保证所有待除的量有限，非零情况单独处理避免除零。$p>1$ 服务于这一 Hölder 论证，$p=1$ 另由逐点三角不等式解决；$p<1$ 的一般障碍已由前例验证。
\end{remark}
% END SUPPLEMENT: id=v14-ch06-minkowski-details
% SOURCE p=198; printed=183; object=ch06-banach-introduction; region=middle
这一结果说明：当 $p\geqslant1$ 时，$L^p$ 是赋范向量空间。更强的结论如下。
% SOURCE p=198; printed=183; object=ch06-lp-completeness; number=6.6; region=middle-lower
\begin{theorem}\label{theorem:ch06-lp-completeness}
当 $1\leqslant p<\infty$ 时，$L^p$ 是 Banach 空间。
\end{theorem}
\begin{proof}
使用\cref{theorem:ch05-banach-series-criterion}。设 $\{f_k\}\subseteq L^p$ 且 $\sum_{k=1}^{\infty}\|f_k\|_p=B<\infty$。令 $G_n=\sum_{k=1}^n|f_k|$、$G=\sum_{k=1}^{\infty}|f_k|$。对所有 $n$，有 $\|G_n\|_p\leqslant\sum_{k=1}^n\|f_k\|_p\leqslant B$；由单调收敛定理，$\int G^p=\lim_n\int G_n^p\leqslant B^p$。故 $G\in L^p$，特别地，$G(x)<\infty$ 几乎处处，这说明级数 $\sum_{k=1}^{\infty}f_k$ 几乎处处收敛。记其和为 $F$，则 $|F|\leqslant G$，所以 $F\in L^p$。又有 $|F-\sum_{k=1}^nf_k|^p\leqslant(2G)^p\in L^1$，故由控制收敛定理，
\begin{align*}
\left\|F-\sum_{k=1}^nf_k\right\|_p^p
=\int\left|F-\sum_{k=1}^nf_k\right|^p\longrightarrow0.
\end{align*}
因此级数 $\sum_{k=1}^{\infty}f_k$ 在 $L^p$ 范数下收敛。
\end{proof}
% SUPPLEMENT: id=v14-ch06-lp-completeness-details; parent=ch06-lp-completeness; kind=proof-detail+motivation+hypotheses; origin=independent-derivation
\begin{proof}
\textbf{代表与极限步骤补写（重构新增）。}先为各 $f_k$ 选可测代表；非负级数 $G$ 可测，且 $G_n^p\uparrow G^p$，因此单调收敛的假设满足。令 $N=\{G=\infty\}$，积分有限给 $\mu(N)=0$。在 $X\setminus N$ 上定义 $F$ 为绝对收敛级数的和，在 $N$ 上置零；这是可测函数，因为可测部分和在可测集 $X\setminus N$ 上逐点收敛。所有控制估计只需在此满测集成立，$2^pG^p$ 是可积控制函数，故所用控制收敛合法。不同代表只会改变一个可数并零集上的值，不改变所得 $L^p$ 元素。范数的 $p$ 次幂趋零与范数趋零等价，因为 $p>0$。先前的级数判据只依赖一般赋范空间论证，不预先使用本定理，因此无循环。
\end{proof}
\begin{remark}
\textbf{动机与逐条件说明（重构新增）。}把每项取绝对值并累加，得到一个同时控制所有部分和的可积包络，从而使抽象完备性归结为两条积分收敛定理。$p\geqslant1$ 使有限部分和满足 Minkowski，并允许使用 Banach 级数判据；有限 $p$ 使积分控制和取 $p$ 次幂有效。测度空间不必完备，因为仅在明确的可测零集上置零；也无有限或 $\sigma$ 有限假设。此定理还补上前章 $L^2$ 为 Hilbert 空间所需的完备性。
\end{remark}
% END SUPPLEMENT: id=v14-ch06-lp-completeness-details
% SOURCE p=198; printed=183; object=ch06-simple-density; number=6.7; region=bottom
\begin{proposition}\label{proposition:ch06-simple-density}
若 $1\leqslant p<\infty$，则所有形如 $f=\sum_{j=1}^na_j\chi_{E_j}$、且每个 $\mu(E_j)<\infty$ 的简单函数，在 $L^p$ 中稠密。
\end{proposition}
\begin{proof}
这些函数显然属于 $L^p$。给定 $f\in L^p$，由\cref{theorem:ch02-simple-approximation}选简单函数序列 $\{f_n\}$，使 $f_n\to f$ 几乎处处且 $|f_n|\leqslant|f|$。于是 $f_n\in L^p$，且 $|f_n-f|^p\leqslant2^p|f|^p\in L^1$。由控制收敛定理，$\|f_n-f\|_p\to0$。若 $f_n=\sum a_j\chi_{E_j}$，其中 $E_j$ 不交且 $a_j\ne0$，则
\begin{align*}
\sum_j|a_j|^p\mu(E_j)=\int|f_n|^p<\infty,
\end{align*}
所以每个 $\mu(E_j)<\infty$。
\end{proof}
% SUPPLEMENT: id=v14-ch06-simple-density-remark; parent=ch06-simple-density; kind=motivation+hypotheses; origin=independent-derivation
\begin{remark}
\textbf{动机与条件说明（重构新增）。}简单函数把一般积分问题降到有限多个集合上的计算。有限测度保证非零常值项可积；在反向验证近似项时，删去系数为零的水平集，再利用每个非零 $|a_j|^p>0$，方可由有限和推出各集测度有限，零水平集本身可以无限。可测性继承自简单逼近，有限 $p$ 使误差有可积控制。原书陈述 $p\geqslant1$ 是为了范数空间背景；同一误差积分推导对 $0<p<1$ 也成立，不过此时稠密性需用后面习题定义的度量理解。无限测度空间上的 $L^\infty$ 不能要求逼近函数都只有有限测度支集，例如常数1到任一此类简单函数的本质上确界距离至少为1。
\end{remark}
% END SUPPLEMENT: id=v14-ch06-simple-density-remark
% SOURCE p=199; printed=184; object=ch06-linfty-introduction; region=top
为补全 $L^p$ 空间的图景，引入对应于极限值 $p=\infty$ 的空间。
% SOURCE p=199; printed=184; object=ch06-essential-supremum-definition; region=upper
\begin{definition}[本质上确界]\label{definition:ch06-essential-supremum}
若 $f$ 是 $X$ 上的可测函数，定义
\begin{align*}
\|f\|_\infty=\inf\{a\geqslant0:\mu(\{x:|f(x)|>a\})=0\},
\end{align*}
约定 $\inf\varnothing=\infty$。$\|f\|_\infty$ 称为 $|f|$ 的\textbf{本质上确界}，有时写成
\begin{align*}
\|f\|_\infty=\operatorname*{ess\,sup}_{x\in X}|f(x)|.
\end{align*}
\end{definition}
% SOURCE p=199; printed=184; object=ch06-essential-bound-attainment; region=upper
这个下确界实际上能够取得，因为
\begin{align*}
\{x:|f(x)|>a\}=\bigcup_{n=1}^{\infty}\{x:|f(x)|>a+n^{-1}\},
\end{align*}
若右边各集合均为零集，左边也是零集。
% SUPPLEMENT: id=v14-ch06-essential-supremum-details; parent=ch06-essential-supremum-definition; kind=proof-detail+motivation+hypotheses; origin=independent-derivation
\begin{proof}
\textbf{有限下确界取得的细节（重构新增）。}令 $A=\{a\geqslant0:\mu(|f|>a)=0\}$。它向上封闭。若 $M=\inf A<\infty$，则 $A\ne\varnothing$；对每个 $n$，由下确界性质可取 $a_n\in A$ 且 $a_n<M+1/n$，因而 $M+1/n\in A$。代入原文的可数并恒等式得到 $\mu(|f|>M)=0$，即 $M\in A$。若 $M=\infty$，则没有有限的本质上界；原文“取得”在这一情形应按扩展上界理解，不能声称空集合内存在有限最小元。
\end{proof}
\begin{remark}
\textbf{动机与条件说明（重构新增）。}本质上确界忽略零集上的尖峰，故适合几乎处处等价类；普通上确界则会被单点改值影响。可测性保证水平集在 $\mathscr M$ 中，使用可数并而非不可数并，使最优本质上界确实成立。取 $a\geqslant0$ 与模长非负相配，空集约定统一处理没有有限上界的情况。这里没有把“上界可取得”误说成“函数在某一点取到最大模长”，后者通常不成立。
\end{remark}
% END SUPPLEMENT: id=v14-ch06-essential-supremum-details
% SOURCE p=199; printed=184; object=ch06-linfty-definition; region=middle
\begin{definition}[$L^\infty$ 空间]\label{definition:ch06-linfty}
定义
\begin{align*}
L^\infty=L^\infty(X,\mathscr M,\mu)=\{f:X\to\mathbb C:f\text{ 可测且 }\|f\|_\infty<\infty\},
\end{align*}
仍约定几乎处处相等的函数确定同一元素。
\end{definition}
% SOURCE p=199; printed=184; object=ch06-linfty-bounded-representative; region=middle
\begin{proposition}\label{proposition:ch06-linfty-bounded-representative}
$f\in L^\infty$ 当且仅当存在有界可测函数 $g$，使 $f=g$ 几乎处处。
\end{proposition}
\begin{proof}
可以取 $g=f\chi_E$，其中 $E=\{x:|f(x)|\leqslant\|f\|_\infty\}$。
\end{proof}
% SUPPLEMENT: id=v14-ch06-linfty-definition-details; parent=ch06-linfty-definition; kind=proof-detail+motivation+hypotheses; origin=independent-derivation
\begin{proof}
\textbf{反向与代表无关性补写（重构新增）。}若 $f\in L^\infty$，刚证明的本质上界取得性质说明 $E^c$ 是可测零集，故原构造的 $g$ 可测、处处有界且与 $f$ 几乎处处相等。反之，若可测 $f=g$ 几乎处处且 $|g|\leqslant C$，则 $\{|f|>C\}$ 包含在二者不等的可测零集中，故 $\|f\|_\infty\leqslant C$。更一般地，若 $f=g$ 几乎处处，则对每个 $a$，两个超水平集只差一个可测零集，于是其本质上界集合相同，$\|f\|_\infty=\|g\|_\infty$。因此定义与代表无关。
\end{proof}
\begin{remark}
\textbf{动机与条件说明（重构新增）。}\cref{definition:ch06-linfty,proposition:ch06-linfty-bounded-representative}把“忽略零集后有统一界”转成“等价类中有处处有界的可测代表”。要求 $g$ 可测，是为了它仍属于同一函数空间；有限本质上界保证截去例外集后真有有限界。无需测度完备，因为使用的 $E$ 及二者不等集合本来可测。只取有界代表并不意味着等价类中每个代表都处处有界。
\end{remark}
% END SUPPLEMENT: id=v14-ch06-linfty-definition-details
% SOURCE p=199; printed=184; object=ch06-linfty-measure-dependence; region=middle
\begin{remark}
两点说明。首先，对固定的 $X,\mathscr M$，$L^\infty(X,\mathscr M,\mu)$ 对 $\mu$ 的依赖，只在于哪些集合具有零测度；若 $\mu,\nu$ 相互绝对连续，则 $L^\infty(\mu)=L^\infty(\nu)$。其次，若 $\mu$ 不是半有限测度，在某些问题中采用稍有不同的 $L^\infty$ 定义更为合适；这一点将在\cref{exercise:ch06-linfty-semifinite-23,exercise:ch06-linfty-semifinite-24,exercise:ch06-linfty-semifinite-25}中探讨。
\end{remark}
% SUPPLEMENT: id=v14-ch06-linfty-null-ideal; parent=ch06-linfty-measure-dependence; kind=proof-detail+motivation+hypotheses; origin=independent-derivation
\begin{proof}
\textbf{同零集结论补写（重构新增）。}相互绝对连续恰说明 $\mu(E)=0$ 与 $\nu(E)=0$ 对每个可测 $E$ 等价。因此两测度对每个 $f$ 给出同一可用本质上界集合，也给出同一几乎处处等价关系；不但空间相同，其本质上确界范数也相同。固定 $\mathscr M$ 保证讨论的是同一批可测代表，单向绝对连续不足以保证两套零集相同。
\end{proof}
% END SUPPLEMENT: id=v14-ch06-linfty-null-ideal
% SOURCE p=199; printed=184; object=ch06-linfty-theorem-introduction; region=middle-lower
此前对 $1\leqslant p<\infty$ 证明的结论，很容易推广到 $p=\infty$，如下所示。
% SOURCE p=199; printed=184; object=ch06-linfty-properties; number=6.8; region=lower
\begin{theorem}\label{theorem:ch06-linfty-properties}
\begin{enumerate}
\item\label{item:ch06-linfty-holder} 若 $f,g$ 为 $X$ 上可测函数，则 $\|fg\|_1\leqslant\|f\|_1\|g\|_\infty$。若 $f\in L^1$ 且 $g\in L^\infty$，则 $\|fg\|_1=\|f\|_1\|g\|_\infty$ 当且仅当在 $f(x)\ne0$ 的集合上，$|g(x)|=\|g\|_\infty$ 几乎处处成立。
\item\label{item:ch06-linfty-norm} $\|\cdot\|_\infty$ 是 $L^\infty$ 上的范数。
\item\label{item:ch06-linfty-uniform} $\|f_n-f\|_\infty\to0$ 当且仅当存在 $E\in\mathscr M$，使 $\mu(E^c)=0$ 且 $f_n\to f$ 在 $E$ 上一致成立。
\item\label{item:ch06-linfty-complete} $L^\infty$ 是 Banach 空间。
\item\label{item:ch06-linfty-simple-density} 简单函数在 $L^\infty$ 中稠密。
\end{enumerate}
\end{theorem}
% SOURCE p=199; printed=184; object=ch06-linfty-properties-proof-omitted; region=lower
证明留给读者（\cref{exercise:ch06-linfty-properties}）。
% SUPPLEMENT: id=v14-ch06-linfty-properties-full-proof; parent=ch06-linfty-properties; kind=proof+motivation+hypotheses; origin=independent-derivation
\begin{remark}\phantomsection\label{remark:ch06-linfty-uniform-scope}
\textbf{集合与量词说明（重构新增）。}第三款的条件是 $\mu(E^c)=0$，即一致收敛发生在一个固定的满测集 $E$ 上。不能漏掉补集上标而换成 $\mu(E)=0$：后一说法取 $E=\varnothing$ 对任意序列都成立，例如在 $[0,1]$ 的 Lebesgue 测度上取 $f_n=1,f=0$，其范数距离恒为1。满测集上的一致收敛，也可等价表述为在某个可测零集之外一致收敛。
\end{remark}
\begin{proof}\phantomsection\label{proof:ch06-linfty-properties-complete}
\textbf{五款完整补证（重构新增）。}
\begin{enumerate}
\item 记 $M=\|g\|_\infty$。若 $\|f\|_1=0$ 或 $M=0$，则乘积几乎处处为零；按非负扩展积分的零乘无穷约定，不等式成立。排除此情形后，若右侧为无穷，结论平凡；否则 $|g|\leqslant M$ 几乎处处，所以积分 $|fg|\leqslant M|f|$ 得结论。若 $f\in L^1,g\in L^\infty$，将可测例外零集上的 $g$ 置零，可令该不等式处处成立。差函数 $|f|(M-|g|)\geqslant0$ 可积，其积分为零当且仅当该函数几乎处处为零；在 $f\ne0$ 的集合上，这恰等价于 $|g|=M$ 几乎处处。修改零集不改变判据，且 $f=0$ 或 $M=0$ 的情况也包括在内。
\item 各函数可取有界可测代表，所以加法和标量乘法封闭并与等价类相容。非负性来自定义；$\|f\|_\infty=0$ 当且仅当 $|f|\leqslant0$ 几乎处处，即 $f$ 是零元素。$c=0$ 时齐次性成立；$c\ne0$ 时 $\{|cf|>a\}=\{|f|>a/|c|\}$，对本质上界取下确界得 $\|cf\|_\infty=|c|\|f\|_\infty$。在两个例外零集的并以外，$|f+g|\leqslant\|f\|_\infty+\|g\|_\infty$，给出三角不等式。
\item 设 $a_n=\|f_n-f\|_\infty\to0$，选 $n_0$ 使 $a_n<\infty$ 对 $n\geqslant n_0$ 成立。各 $N_n=\{|f_n-f|>a_n\}$ 为可测零集，故 $N=\bigcup_{n\geqslant n_0}N_n$ 仍是可测零集。令 $E=X\setminus N$，则 $\mu(E^c)=0$，且在 $E$ 上，$\sup|f_n-f|\leqslant a_n\to0$，所以一致收敛。反之，若在 $E$ 上一致收敛且 $\mu(E^c)=0$，则对每个 $\varepsilon>0$，充分大 $n$ 满足 $|f_n-f|\leqslant\varepsilon$ 于 $E$，故 $\|f_n-f\|_\infty\leqslant\varepsilon$。这也表明例外零集必须对整个尾序列统一，不能每个 $n$ 随意忽略不同且未合并的集合。
\item 设 $(f_n)$ 在本质上确界范数中 Cauchy，先选每个 $f_n$ 的有界可测代表。对所有 $(n,m)\in\mathbb N^2$，删去可测零集 $\{|f_n-f_m|>\|f_n-f_m\|_\infty\}$ 的可数并 $N$。在 $X\setminus N$ 上，$(f_n)$ 一致 Cauchy，因此由 $\mathbb C$ 完备逐点定义极限 $f$。给定 $\varepsilon>0$，取 $K$ 使 $n,m\geqslant K$ 时范数距离小于 $\varepsilon$，固定 $n\geqslant K$ 后令 $m\to\infty$，得 $|f_n-f|\leqslant\varepsilon$ 于 $X\setminus N$，故一致收敛。在 $N$ 上将 $f$ 置零，它可测；用 $\varepsilon=1$ 和某个有界 $f_K$ 得 $|f|\leqslant\|f_K\|_u+1$，所以 $f\in L^\infty$，且刚才估计给范数收敛。
\item 给定 $f\in L^\infty$，取处处有界可测代表 $g$。对每个 $n\geqslant1$，把包含 $g(X)$ 的复平面有界矩形分成有限个边长不超过 $1/n$ 的半开小矩形，并妥善包含外边界。在各小矩形原像上，取该矩形一个固定点作为值，得可测简单函数 $s_n$；所有点的误差至多 $\sqrt2/n$。于是 $\|f-s_n\|_\infty\leqslant\sqrt2/n\to0$。这里小矩形原像不必具有有限测度。
\end{enumerate}
\end{proof}
\begin{remark}
\textbf{动机与逐条件说明（重构新增）。}第一款用本质上界直接控制乘积；等号要求在 $f$ 真正参与积分的地方，用尽 $g$ 的最大允许模长，所以无须限制 $f=0$ 处的 $g$。第二款的等价类识别负责范数正定性；第三款的可数序列允许把全部例外集合合成一个零集，结论不是单纯逐点几乎处处收敛。第四款依赖标量完备性以及一致 Cauchy 估计；第五款依赖有限本质上界，使值域只需有限网格。所有五款只要求测度空间，不要求有限、半有限、$\sigma$ 有限或完备测度；每次修改的集合都可测。第五款与有限 $p$ 的稠密命题不同，并不要求简单函数的非零水平集测度有限。
\end{remark}
% END SUPPLEMENT: id=v14-ch06-linfty-properties-full-proof
% SOURCE p=199; printed=184; object=ch06-endpoint-conjugate-definition; region=bottom
\begin{definition}[端点共轭指数]\label{definition:ch06-endpoint-conjugate}
根据\subcref{theorem:ch06-linfty-properties}{item:ch06-linfty-holder}以及形式等式 $1^{-1}+\infty^{-1}=1$，自然地把 $1$ 与 $\infty$ 视为互为共轭指数；此后采用这一约定。
\end{definition}
% SUPPLEMENT: id=v14-ch06-endpoint-conjugate-remark; parent=ch06-endpoint-conjugate-definition; kind=motivation+hypotheses; origin=independent-derivation
\begin{remark}
\textbf{动机与约定说明（重构新增）。}这里的 $\infty^{-1}=0$ 是扩展指数记号，使有限共轭公式连续衔接端点；对应的估计确由上一条端点 Hölder 证明支持。它不是把无穷当作可任意作实数运算的数，也不将有限指数证明中的除以 $p-1$ 操作套到 $p=1$。
\end{remark}
% END SUPPLEMENT: id=v14-ch06-endpoint-conjugate-remark
% SOURCE p=199,200; printed=184,185; object=ch06-uniform-essential-comparison; region=bottom-top
\subcref{theorem:ch06-linfty-properties}{item:ch06-linfty-uniform}表明，$\|\cdot\|_\infty$ 与一致范数 $\|\cdot\|_u$ 密切相关，但通常不同。
% SOURCE p=199,200; printed=184,185; object=ch06-continuous-essential-supremum; region=bottom-top
\begin{proposition}\label{proposition:ch06-continuous-essential-supremum}
对于 Lebesgue 测度，或更一般地，对每个开集赋予正值的 Borel 测度，只要 $f$ 连续，就有 $\|f\|_\infty=\|f\|_u$。此时可交替使用两种记号，并把有界连续函数空间视为 $L^\infty$ 的一个闭子空间。
\end{proposition}
\begin{proof}
集合 $\{x:|f(x)|>a\}$ 是开集。
\end{proof}
% SUPPLEMENT: id=v14-ch06-continuous-essential-details; parent=ch06-continuous-essential-supremum; kind=proof+source-clarification+motivation+hypotheses; origin=independent-derivation
\begin{remark}
\textbf{原文范围说明（重构新增）。}“每个开集有正测度”应指每个\textbf{非空}开集，因为空集总有零测度。以下采用这一准确含义；允许函数的一致上确界为无穷。
\end{remark}
\begin{proof}
\textbf{范数相等及闭性补证（重构新增）。}总有 $\|f\|_\infty\leqslant\|f\|_u$。若 $0\leqslant a<\|f\|_u$，上确界定义保证存在 $x$ 使 $|f(x)|>a$，于是相应超水平集为非空开集，具有正测度。这说明 $a$ 不是本质上界，所以 $\|f\|_\infty\geqslant\|f\|_u$；若后一数为无穷，对任意有限 $a$ 应用同一推论。零上确界时等式直接成立。

把 $g\in C_b(X)$ 送到其几乎处处等价类。范数相等给线性等距性，并给单射性。若这一子空间中的序列 $[g_n]$ 在 $L^\infty$ 中趋于 $h$，则 $(g_n)$ 在一致范数中 Cauchy。逐点用 $\mathbb C$ 完备性取极限 $g$，一致 Cauchy 估计使收敛一致；因此 $g$ 连续，且由某个 $g_N$ 的有界性和一致误差控制得到 $g$ 有界。于是 $[g_n]\to[g]$ 于 $L^\infty$，极限唯一给 $h=[g]$，所以子空间序列闭。范数空间可度量，序列闭即闭：若一点在闭包中，每个半径 $1/n$ 球取一点即可得到逼近序列。
\end{proof}
\begin{remark}
\textbf{动机与条件说明（重构新增）。}连续性把某点处超出阈值的现象扩展到开邻域；非空开集正测度则阻止这种现象被零集等价关系抹去。Borel 测度保证这些开集可测。两条件共同让点态连续函数能无歧义地嵌入等价类空间。有界性用于进入 $L^\infty$，而范数相等本身可允许无穷。若去掉开集正测度这一背景，结论一般失败：在 $\mathbb R$ 上取集中于0的 Dirac 测度，连续有界函数 $g(x)=\min(1,|x|)$ 的本质上确界为0，一致上确界为1。这里不需要 $X$ 紧或测度有限。
\end{remark}
% END SUPPLEMENT: id=v14-ch06-continuous-essential-details
% SOURCE p=200; printed=185; object=ch06-lp-inclusions-introduction; region=upper
一般而言，当 $p\ne q$ 时，$L^p$ 并不包含于 $L^q$。为理解问题所在，考虑 $(0,\infty)$ 上 Lebesgue 测度的以下简单例子。
% SOURCE p=200; printed=185; object=ch06-power-examples; region=upper
\begin{examples}\label{examples:ch06-power-integrability}
令 $f_a(x)=x^{-a}$，其中 $a>0$。
\begin{itemize}
\item 初等微积分表明，$f_a\chi_{(0,1)}\in L^p$ 当且仅当 $p<a^{-1}$。
\item $f_a\chi_{(1,\infty)}\in L^p$ 当且仅当 $p>a^{-1}$。
\end{itemize}
\end{examples}
% SUPPLEMENT: id=v14-ch06-power-examples-details; parent=ch06-power-examples; kind=proof+motivation+hypotheses; origin=independent-derivation
\begin{proof}
\textbf{两端积分补写（重构新增）。}对有限 $p>0$，令 $b=ap>0$。若 $b\ne1$，
\begin{align*}
\int_\delta^1x^{-b}\,\mathrm dx=\frac{1-\delta^{1-b}}{1-b},
\qquad
\int_1^Rx^{-b}\,\mathrm dx=\frac{R^{1-b}-1}{1-b}.
\end{align*}
令 $\delta\downarrow0$ 或 $R\uparrow\infty$，前者有限当且仅当 $b<1$，后者有限当且仅当 $b>1$。$b=1$ 时分别为 $-\ln\delta$ 和 $\ln R$，均发散。对 $p=\infty$，前一函数无本质上界，因为任意高阈值仍在一段正长度区间上被超过；后一函数的本质上确界为1。这与所写严格指数条件一致。
\end{proof}
\begin{remark}
\textbf{动机与条件说明（重构新增）。}当 $p<q$，在原点取 $1/q\leqslant a<1/p$，便得属于 $L^p$ 而不属于 $L^q$ 的第一类函数；在无穷远取 $1/q<a\leqslant1/p$，便得属于 $L^q$ 而不属于 $L^p$ 的第二类函数。这里先取有限 $q$；若 $q=\infty$，第一类取 $0<a<1/p$，第二类取 $a=1/p$ 即可。正参数使原点成为奇性而无穷远成为衰减端，Lebesgue 测度把二者转成明确瑕积分。原文“一般而言”不是对每个固定测度空间否认包含关系；随后两个命题正提供例外背景。
\end{remark}
% END SUPPLEMENT: id=v14-ch06-power-examples-details
% SOURCE p=200; printed=185; object=ch06-local-global-discussion; region=upper-middle
由此看到函数 $f$ 不属于 $L^p$ 的两种原因：$|f|^p$ 可能在某点附近增长得过快，也可能在无穷远衰减得不够快。前一种情形中，$p$ 增大使行为更坏；后一种情形则更好。换言之，若 $p<q$，$L^p$ 中的函数局部可以比 $L^q$ 中的函数更加奇异，而 $L^q$ 中的函数整体上可以比 $L^p$ 中的函数分布得更散。这些略显粗略的想法，实际上能相当准确地指引一般情形。下面给出四个精确结果；后两个说明，若测度空间的条件排除了上述某一种坏行为，就能得到 $L^p\subseteq L^q$ 一类的包含关系。更一般的结果见\cref{exercise:ch06-lp-inclusion-criterion}。
% SOURCE p=200; printed=185; object=ch06-lp-sum-inclusion; number=6.9; region=middle
\begin{proposition}\label{proposition:ch06-lp-sum-inclusion}
若 $0<p<q<r\leqslant\infty$，则 $L^q\subseteq L^p+L^r$；即每个 $L^q$ 函数都是一个 $L^p$ 函数和一个 $L^r$ 函数之和。
\end{proposition}
\begin{proof}
若 $f\in L^q$，令 $E=\{x:|f(x)|>1\}$，并置 $g=f\chi_E$、$h=f\chi_{E^c}$。则 $|g|^p=|f|^p\chi_E\leqslant|f|^q\chi_E$，所以 $g\in L^p$；又 $|h|^r=|f|^r\chi_{E^c}\leqslant|f|^q\chi_{E^c}$，所以 $h\in L^r$。（$r=\infty$ 时，显然 $\|h\|_\infty\leqslant1$。）
\end{proof}
% SUPPLEMENT: id=v14-ch06-sum-inclusion-remark; parent=ch06-lp-sum-inclusion; kind=motivation+hypotheses; origin=independent-derivation
\begin{remark}
\textbf{动机与条件说明（重构新增）。}以模长1为界，大值部分在较小指数下更容易积分，小值部分在较大指数下更容易积分。$p<q<r$ 正是这两项点态幂比较所用的顺序；$p>0$ 保证幂仍为非负实量，$q<\infty$ 给出控制积分。$E$ 可测来自 $f$ 可测，$g+h=f$ 处处成立；$r=\infty$ 必须用有界性单独论证，不写 $|h|^\infty$。截断阈值1仅为方便，任意正阈值都可并产生相应常数；这不是固定阈值的必要性主张。
\end{remark}
% END SUPPLEMENT: id=v14-ch06-sum-inclusion-remark
% SOURCE p=200; printed=185; object=ch06-lp-interpolation; number=6.10; region=lower
\begin{proposition}\label{proposition:ch06-lp-interpolation}
若 $0<p<q<r\leqslant\infty$，则 $L^p\cap L^r\subseteq L^q$，且
\begin{align*}
\|f\|_q\leqslant\|f\|_p^\lambda\|f\|_r^{1-\lambda},
\end{align*}
其中 $\lambda\in(0,1)$ 由下式定义：
\begin{align*}
q^{-1}=\lambda p^{-1}+(1-\lambda)r^{-1},
\qquad
\lambda=\frac{q^{-1}-r^{-1}}{p^{-1}-r^{-1}}.
\end{align*}
\end{proposition}
\begin{proof}
若 $r=\infty$，则 $|f|^q\leqslant\|f\|_\infty^{q-p}|f|^p$，且 $\lambda=p/q$，所以
\begin{align*}
\|f\|_q\leqslant\|f\|_p^{p/q}\|f\|_\infty^{1-p/q}
=\|f\|_p^\lambda\|f\|_\infty^{1-\lambda}.
\end{align*}
若 $r<\infty$，应用 Hölder 不等式，共轭指数取 $p/(\lambda q)$ 与 $r/((1-\lambda)q)$：
\begin{align*}
\int|f|^q
&=\int|f|^{\lambda q}|f|^{(1-\lambda)q}
\leqslant\bigl\||f|^{\lambda q}\bigr\|_{p/(\lambda q)}
       \bigl\||f|^{(1-\lambda)q}\bigr\|_{r/((1-\lambda)q)}\\
&=\left[\int|f|^p\right]^{\lambda q/p}
  \left[\int|f|^r\right]^{(1-\lambda)q/r}
=\|f\|_p^{\lambda q}\|f\|_r^{(1-\lambda)q}.
\end{align*}
取 $q$ 次方根即得结论。
\end{proof}
% SUPPLEMENT: id=v14-ch06-interpolation-details; parent=ch06-lp-interpolation; kind=proof-detail+motivation+hypotheses; origin=independent-derivation
\begin{remark}
\textbf{动机与逐条件说明（重构新增）。}中间倒数指数是两端倒数的严格凸组合，所以 $0<\lambda<1$。有限 $r$ 时，两数 $\lambda q/p$ 与 $(1-\lambda)q/r$ 都正且和为1，因而各自小于1；这核实了实际用于 Hölder 的指数均严格大于1，即使原来的 $p<1$ 也完全合法。两端的有限大小保证右端有限，若任一为零，则 $f=0$ 几乎处处，结论另行直接成立。$r=\infty$ 的点态估计按几乎处处理解，零范数情况同样先处理。插值在这里通过两个端点控制中间大小，无须任何有限测度假设。
\end{remark}
% END SUPPLEMENT: id=v14-ch06-interpolation-details
% SOURCE p=201; printed=186; object=ch06-sequence-inclusion; number=6.11; region=top
\begin{proposition}\label{proposition:ch06-sequence-inclusion}
若 $A$ 为任意集合，且 $0<p<q\leqslant\infty$，则 $\ell^p(A)\subseteq\ell^q(A)$，且 $\|f\|_q\leqslant\|f\|_p$。
\end{proposition}
\begin{proof}
显然 $\|f\|_\infty^p=\sup_\alpha|f(\alpha)|^p\leqslant\sum_\alpha|f(\alpha)|^p$，所以 $\|f\|_\infty\leqslant\|f\|_p$。$q<\infty$ 的情形由\cref{proposition:ch06-lp-interpolation}得到：取 $\lambda=p/q$，则
\begin{align*}
\|f\|_q\leqslant\|f\|_p^\lambda\|f\|_\infty^{1-\lambda}\leqslant\|f\|_p.
\end{align*}
\end{proof}
% SUPPLEMENT: id=v14-ch06-sequence-inclusion-remark; parent=ch06-sequence-inclusion; kind=motivation+hypotheses; origin=independent-derivation
\begin{remark}
\textbf{动机与条件说明（重构新增）。}计数测度中单个非零坐标所在的单点已有测度1，不能把很高的值藏在任意小的正测度集合里，因此有限 $p$ 大小自动控制每个坐标。任意指标集的非负和仍支配每项，不要求 $A$ 可数。有限 $q$ 用 $p<q<\infty$ 的插值，$q=\infty$ 已由逐坐标控制直接完成；$A=\varnothing$ 时所有函数空间均为零空间，结论按零范数解释。若计数测度改为具有任意小正权重的离散测度，常数1的这一步估计便不能照搬。
\end{remark}
% END SUPPLEMENT: id=v14-ch06-sequence-inclusion-remark
% SOURCE p=201; printed=186; object=ch06-finite-measure-inclusion; number=6.12; region=middle
\begin{proposition}\label{proposition:ch06-finite-measure-inclusion}
若 $\mu(X)<\infty$ 且 $0<p<q\leqslant\infty$，则 $L^p(\mu)\supseteq L^q(\mu)$，且
\begin{align*}
\|f\|_p\leqslant\|f\|_q\mu(X)^{1/p-1/q}.
\end{align*}
\end{proposition}
\begin{proof}
若 $q=\infty$，这是显然的：
\begin{align*}
\|f\|_p^p=\int|f|^p\leqslant\|f\|_\infty^p\int1
=\|f\|_\infty^p\mu(X).
\end{align*}
若 $q<\infty$，以 $q/p$ 与 $q/(q-p)$ 为共轭指数应用 Hölder 不等式：
\begin{align*}
\|f\|_p^p=\int|f|^p\cdot1
\leqslant\bigl\||f|^p\bigr\|_{q/p}\|1\|_{q/(q-p)}
=\|f\|_q^p\mu(X)^{(q-p)/q}.
\end{align*}
\end{proof}
% SUPPLEMENT: id=v14-ch06-finite-measure-inclusion-remark; parent=ch06-finite-measure-inclusion; kind=proof-detail+motivation+hypotheses; origin=independent-derivation
\begin{remark}
\textbf{动机与条件说明（重构新增）。}有限总测度排除了函数把很小的值铺展在无穷大区域上所造成的不可积性。$q>p>0$ 使实际 Hölder 指数 $q/p$、$q/(q-p)$ 大于1并共轭；总测度有限正是常数函数1属于第二个空间的条件。取原估计的 $p$ 次方根得到所述指数差。若 $\mu(X)=0$，所有等价类都为零，指数差为正，原式两边均为零；无需除以总测度。$q=\infty$ 时约定 $1/q=0$ 并用单独的逐点估计。删除有限性不能得到对一般空间的同一包含结论：在 $(1,\infty)$ 的 Lebesgue 测度上，前述 $x^{-1/p}$ 属于每个有限 $q>p$ 及 $L^\infty$，却不属于 $L^p$。这只说明该一般结论不能无条件推广，并不声称有限总测度是所有特定空间发生包含的必要条件。
\end{remark}
% END SUPPLEMENT: id=v14-ch06-finite-measure-inclusion-remark
% SOURCE p=201; printed=186; object=ch06-lp-significance; region=middle-lower
最后谈几句 $L^p$ 空间的意义。最显然重要的三个是 $L^1$、$L^2$ 与 $L^\infty$。我们已熟悉 $L^1$；$L^2$ 特别之处在于它是 Hilbert 空间；$L^\infty$ 上的拓扑则与一致收敛拓扑密切相关。遗憾的是，$L^1$ 与 $L^\infty$ 在许多方面有病态表现，研究中间的 $L^p$ 空间往往更有成果。\cref{section:ch06-lp-dual}中的对偶理论是一个体现；另一个体现是：Fourier 分析与微分方程中许多重要算子在 $1<p<\infty$ 的 $L^p$ 上有界，却不在 $L^1$ 或 $L^\infty$ 上有界。（\cref{section:ch09-fourier-applications}将提到一些例子。）
\subfile{section-01-exercises.tex}
\end{document}
